\scanpage{37}
\begingroup\makeatletter\tagsleft@true\makeatother
\uline{\unsure{Cor|Cor.}} Dans le cas où $\operatorname{Card}\{\ \}=4$, \inserted{et $\Sigma_1^-\ne X,\Sigma_2^-\ne X$} les seules relations d'inclusions entre les ensembles $\Sigma_1^+,\Sigma_1^-,\Sigma_2^+,\Sigma_2^-$ sont les relations
% slanted
\margin{Sous forme : Pour que les relations \unsure{strictes} \unsure{(*)} soient les seules relations d'inclusion entre les ensembles $\Sigma_1^+,\Sigma_1^-,\Sigma_2^+,\Sigma_{\unsure{2|1}}^-$, il \unsure{f|f.} et s. que ceux-ci soient \inserted{\unsure{mutlt.}} distincts, et que l'on ait $\Sigma_1^-\ne X,\ \Sigma_2^-\ne X$.}
\[
\left\{\begin{aligned}
\Sigma_1^+&\subset\Sigma_2^-\\
\Sigma_2^+&\subset\Sigma_1^-
\end{aligned}\right. \quad .
\tag*{($*$)}
\]
En effet, on n'a pas les inclusions inverses, par hypothèse. D'autre part, on n'\uline{a pas}
\[
\Sigma_1^+\subset\Sigma_1^- \quad\text{ou}\quad \Sigma_1^-\unsure{\subset|\in}\Sigma_1^+,\quad\text{car on aurait}
\]
$\deleted{\unsure{\Sigma_1^+,\Sigma_1^+}}\ \Sigma_1^-=X$, de $=$ on n'a pas $\Sigma_2^+\subset{\unsure{\Sigma|E}}_{\unsure{2|1}}^-$ ni $\Sigma_2^-\subset\Sigma_2^+$, et on n'a pas $\Sigma_1^+\subset\Sigma_2^+$ ni l'inclusion inverse, puisque $\Sigma_1^+$ et $\unsure{\Sigma^+|\Sigma_2^+}$ sont disjoints et $\ne\varnothing$. Il reste à voir que l'on n'a pas une inclusion $\Sigma_1^-\subset\Sigma_2^-$ ou inclusion inverse, on lit sur le diagramme $(14)$ que \deleted{\unsure{cela}} \inserted{\unsure{la première|le premier} \unsure{p.ex.|p. ex.}} signifierait que $\deleted{\Sigma_{12}\unsure{\simeq|=}\Sigma_1^-}$ \deleted{\unsure{d'où}} $\Sigma_2^-=X$, ce que l'on a exclu.

% Vertical rule down the left of the definition.
\noindent\makebox[0pt][r]{\rule[-9ex]{.4pt}{10.5ex}\hspace{.55em}}\uline{\unsure{Déf|Déf.}} On dit que $\Pi_1,\Pi_2$ sont \uline{strictement non panachées} \inserted{si elles sont non panachées, et} si \deleted{\unsure{les}}les quatre ensembles $\Sigma_1^+,\Sigma_1^-,\Sigma_2^+,\Sigma_2^-$ sont distincts, \unsure{et.|i.e.} si \unsure{((|c} de plus $\Sigma_1^-\ne X,\Sigma_2^-\ne X$ -- i.e si les seules relations d'inclusion entre les ensembles $\Sigma_1^+,\Sigma_1^-,\Sigma_2^+,\Sigma_2^-$ sont $\Sigma_1^+\subset\Sigma_2^-,\ \Sigma_2^+\subset\Sigma_1^-$ -- ce qui signifie aussi que l'on a :
\[
\Sigma_1^-\ne X,\Sigma_2^-\ne X
\quad(\text{on dit \unsure{donc|alors} que }\Pi_1,\Pi_2\text{ sont \uline{propres}}),
\ E_1\ne\Sigma_{12},E_2\ne\Sigma_{12}.
\tag*{(21)}
\]

\endgroup
\clearpage
% archive sheet number: 36
